\section{Electroweak embedding as a projective datum}
\label{sec:ew-embedding}
% Status: cited/open. Projective interpretation follows from the electroweak mass matrix; ordered sampling remains open.

The gauge-Higgs sector gives the tree-level vector masses from the Higgs kinetic term
\begin{equation}
  (D_\mu H)^\dagger(D^\mu H),
  \qquad
  D_\mu=\partial_\mu-\frac{i g}{2}\sigma^a W^a_\mu-\frac{i g'}{2}B_\mu ,
  \qquad
  H=\frac{1}{\sqrt2}\begin{pmatrix}0\\ v+h\end{pmatrix}.
  \label{eq:higgs-kinetic}
\end{equation}
The charged fields
\begin{equation}
  W^\pm_\mu=\frac{1}{\sqrt2}\left(W^1_\mu\mp iW^2_\mu\right)
  \label{eq:wpm-fields}
\end{equation}
receive a mass term \(g^2v^2 W^+_\mu W^{-\mu}/4\).  The neutral fields have mass matrix
\begin{equation}
  \frac{v^2}{8}
  \begin{pmatrix} W^3_\mu & B_\mu \end{pmatrix}
  \begin{pmatrix}
    g^2 & -gg'\\
    -gg' & g'^2
  \end{pmatrix}
  \begin{pmatrix} W^{3\mu}\\ B^\mu \end{pmatrix}.
  \label{eq:neutral-mass-matrix}
\end{equation}
It is diagonalized by
\begin{equation}
  Z_\mu=\cos\theta_W\,W^3_\mu-\sin\theta_W\,B_\mu,
  \qquad
  A_\mu=\sin\theta_W\,W^3_\mu+\cos\theta_W\,B_\mu,
  \qquad
  \tan\theta_W=\frac{g'}{g}.
  \label{eq:weak-rotation}
\end{equation}
This gives
\begin{equation}
  M_W^2=\frac{g^2v^2}{4},
  \qquad
  M_Z^2=\frac{(g^2+g'^2)v^2}{4},
  \qquad
  M_\gamma^2=0.
  \label{eq:tree-masses}
\end{equation}
The ratio is
\begin{equation}
  1-\frac{M_W^2}{M_Z^2}=\frac{g'^2}{g^2+g'^2}.
  \label{eq:tree-angle}
\end{equation}
Thus the vacuum scale \(v\) fixes the radial size of the massive vector pair, while the quotient fixes a projective direction in the two-coupling plane \cite{Logan2022,AlvarezGaumeVazquezMozo2023}.

\subsection{Custodial structure and the electroweak ray}

The Standard Model doublet can be packaged as the bidoublet
\begin{equation}
  \Phi_H=\begin{pmatrix}
  H^{0*} & H^+\\
  -H^{+*} & H^0
  \end{pmatrix}.
  \label{eq:higgs-bidoublet}
\end{equation}
In the scalar sector this makes the approximate global \(SU(2)_L\times SU(2)_R\) structure visible; the vacuum proportional to the identity preserves the diagonal custodial \(SU(2)\) subgroup \cite{Logan2022}.  Hypercharge gauging and Yukawa splittings break the symmetry, while the tree-level doublet mass matrix preserves the relation
\begin{equation}
  \rho=\frac{M_W^2}{M_Z^2\cos^2\theta_W}=1.
  \label{eq:rho-tree}
\end{equation}
This is the electroweak reason the DeVries comparison is read as a projective statement about the W/Z pair.  The broken-to-unbroken deformation scales the vector masses with \(v\) while preserving the coupling-space direction.

The DeVries proposal is read in the same projective language:
\begin{equation}
  \frac{M_W^2}{M_Z^2}\quad \longleftrightarrow\quad \frac{x_+(3/4)}{x_+(2)}.
  \label{eq:projective-identification}
\end{equation}
At this stage Eq.~\eqref{eq:projective-identification} is a tree-level
projective analogy.  Sec.~\ref{sec:pole} and
Appendix~\ref{app:target-pole-placement} state the additional pole-matching
obligation.  The common scale is then supplied by the electroweak vacuum or, in
a deeper model, by the compactification or brane scale that generates the
effective gauge-Higgs sector.

The unbroken limit is the ray
\begin{equation}
  v\to0
  \quad\text{with}\quad
  \frac{g'^2}{g^2+g'^2}\;\text{fixed}.
  \label{eq:ew-ray}
\end{equation}
Along this ray, \(M_W\) and \(M_Z\) vanish together and the photon remains massless.  This is the electroweak deformation relevant to the construction.  The projective datum remains meaningful as the limiting direction even when the radial mass scale is taken to zero.

Equivalently, introduce a parameter \(t\) by
\begin{equation}
  H_t=\frac{1}{\sqrt2}
  \begin{pmatrix}0\\t v+h\end{pmatrix},
  \qquad
  g(t)=g,\qquad g'(t)=g',
  \qquad 0\le t\le1 .
  \label{eq:ew-ray-family}
\end{equation}
Then
\begin{equation}
  M_W^2(t)=t^2\frac{g^2v^2}{4},
  \qquad
  M_Z^2(t)=t^2\frac{(g^2+g'^2)v^2}{4},
  \qquad
  M_\gamma^2(t)=0 .
  \label{eq:ew-ray-family-masses}
\end{equation}
This is a tree-level ray statement before the pole-matching map of
Sec.~\ref{sec:pole}.  Coupling-space paths such as \(g\to0\), \(g'\to0\),
\(g'/g\to0\), or \(g'/g\to\infty\), and custodial-breaking shifts such as
\(\delta M_W^2\ne\delta M_Z^2\cos^2\theta_W\), define separate deformation
problems.  A source route for the DeVries determinant must preserve the
single radial parameter, the photon null direction, and the projective W/Z
direction.  Routes with different behavior enter the manuscript as failure
modes or as separate models.

The central analytical question is the ordered assignment
\begin{equation}
  (J_W,J_Z)=\left(\frac34,2\right).
  \label{eq:Jassignment}
\end{equation}
The electroweak gauge fields are gauge-algebra fields before symmetry breaking: \(W^1,W^2,W^3\) live in the adjoint of \(SU(2)_L\), while \(B\) is the gauge field of the abelian \(U(1)_Y\) factor.  The value \(J=2\) has a direct \(SU(2)_L\) adjoint reading through \(T(T+1)\) with \(T=1\).  The value \(J=3/4\) belongs to the Higgs/order-parameter side, where the symmetry-breaking field is an \(SU(2)_L\) doublet with \(T=1/2\) and \(T(T+1)=3/4\).

These representation labels define the available samples:
\begin{equation}
  J_Z=T_{\rm adj}(T_{\rm adj}+1)=2,
  \qquad
  J_H=T_H(T_H+1)=\frac34 .
  \label{eq:adjoint-doublet-casimirs}
\end{equation}
The remaining calculation is the ordered sampling rule that turns the pair \((J_H,J_{\rm adj})\) into the W/Z quotient in Eq.~\eqref{eq:projective-identification}.  A completed derivation must start from the gauge-Higgs Lagrangian, a current-algebra construction, or a compactification boundary problem and produce Eq.~\eqref{eq:Jassignment} from the same determinant that supplies the negative branch.

\subsection{The assignment problem as a theorem target}

The theorem target can be stated mechanism-independently.  Let \(\mathcal E\) be the electroweak gauge-Higgs system, including its projective W/Z pole data and scalar order-parameter channel.  A DeVries derivation is a map
\begin{equation}
  \mathcal D:\mathcal E\longrightarrow \mathcal M_J
  =
  \begin{pmatrix}
  0 & \sqrt J\\
  \sqrt J & -J
  \end{pmatrix}
  \label{eq:assignment-map}
\end{equation}
with the following outputs:
\begin{align}
  \mathcal D(\text{order-parameter channel})&\mapsto J_H=3/4,\\
  \mathcal D(\text{adjoint/current channel})&\mapsto J_{\rm adj}=2,\\
  \frac{M_{W,\rm pole}^2}{M_{Z,\rm pole}^2}
  &=\frac{x_+(J_H)}{x_+(J_{\rm adj})}+\Delta_{\rm match}.
  \label{eq:assignment-outputs}
\end{align}
The exact pole claim requires a derived vanishing of \(\Delta_{\rm match}\) or
a computed matching remainder with stated scheme and source-theory origin.  The
theorem must also account for the negative branch in the same basis, so that
the vector ratio and scalar clue arise from one operator.

A natural current to test is the Higgs isospin current
\begin{equation}
  \mathcal J^a_\mu
  =
  i H^\dagger\frac{\sigma^a}{2}D_\mu H
  -i(D_\mu H)^\dagger\frac{\sigma^a}{2}H,
  \label{eq:higgs-isospin-current}
\end{equation}
which transforms as an adjoint object while being built from the doublet order parameter.  The current contains the two data needed for the ordered pair: doublet data enter the source, adjoint data enter the current.  Turning this observation into Eq.~\eqref{eq:assignment-outputs} is the remaining electroweak theorem target.
